% GENERATED FILE. Edit canonical lessons, then run npm run generate:books.
% canonical-source-hash: fnv1a64:312bcb6505202803
% canonical-lessons: MW-C71-bakri, MW-C71-bher, MW-C71-chiri, MW-C71-nahar, MW-C71-hathi

\chapter{Goat, Sheep, Sparrow, Lion, Elephant}
\label{ch:mw-goat-sheep-sparrow-lion-elephant}
\hlchaptermodality{71}

\begin{chapteropening}
\textbf{By the end of this chapter:} \emph{I can say a goat, a sheep, a sparrow, a small bird, a lion, a tiger and an elephant.}

Complete the last lesson of chapter 71: I can say a goat, a sheep, a sparrow, a small bird, a lion, a tiger and an elephant.
\end{chapteropening}

\section[\texorpdfstring{\mw{बकरी} (bakrī)}{बकरी (bakrī)}]{\mw{बकरी} (bakrī) — a goat}

\label{lesson:MW-C71-bakri}



\begin{quote}
Before the new one: say the Marwadi for a cow, then the Marwadi for a buffalo.
\end{quote}

\subsection*{You'll want to know: \mw{बकरी}}
\textbf{\mw{बकरी}} — \emph{bakrī} — \textquotedblleft{}a goat\textquotedblright{}.

\begin{cousinweb}
\textbf{\mw{बकरी}} (\emph{bakrī}) is a tadbhava, a word inherited from Sanskrit and reshaped by sound change. It comes from \emph{barkara}, \textquotedblleft{}kid, lamb; goat\textquotedblright{}. CDIAL gives no Marwari form in this entry (CDIAL 9153); the match is through Hindi \emph{bakrī}.
\end{cousinweb}

\begin{grammarlens}[title={What you've built}]
One more word for this chapter.
\end{grammarlens}

\subsection*{Your turn}
\begin{itemize}
\raggedright
  \item \emph{Say it:} \emph{bakrī}
  \item \emph{Say it:} \emph{bakrī}, once more
  \item \emph{Say it:} say \emph{bakrī}
  \item \emph{Recall:} say the Marwadi for a cow, then the Marwadi for a buffalo, then say \emph{bakrī} again
\end{itemize}

\begin{tcolorbox}[breakable,colback=teal!4,colframe=teal!35!black,title={Before you move on}]
What does \mw{बकरी} mean? (\textquotedblleft{}A goat\textquotedblright{}.) Say it once more.
\end{tcolorbox}
\section[\texorpdfstring{\mw{भेड़} (bheṛ)}{भेड़ (bheṛ)}]{\mw{भेड़} (bheṛ) — a sheep}

\label{lesson:MW-C71-bher}



\begin{quote}
Before the new one: say the Marwadi for a buffalo, then the Marwadi for a goat.
\end{quote}

\subsection*{You'll want to know: \mw{भेड़}}
\textbf{\mw{भेड़}} — \emph{bheṛ} — \textquotedblleft{}a sheep\textquotedblright{}.

\begin{cousinweb}
\textbf{\mw{भेड़}} (\emph{bheṛ}) is a tadbhava, a word inherited from Sanskrit and reshaped by sound change. It comes from \emph{bheḍra}, \textquotedblleft{}ram\textquotedblright{}. CDIAL gives no Marwari form in this entry (CDIAL 9606); the match is through Hindi \emph{bheṛ}. CDIAL takes the Sanskrit word itself from an Austroasiatic language.
\end{cousinweb}

\begin{grammarlens}[title={What you've built}]
One more word for this chapter.
\end{grammarlens}

\subsection*{Your turn}
\begin{itemize}
\raggedright
  \item \emph{Say it:} \emph{bheṛ}
  \item \emph{Say it:} \emph{bheṛ}, once more
  \item \emph{Say it:} say \emph{bheṛ}
  \item \emph{Recall:} say the Marwadi for a buffalo, then the Marwadi for a goat, then say \emph{bheṛ} again
\end{itemize}

\begin{tcolorbox}[breakable,colback=teal!4,colframe=teal!35!black,title={Before you move on}]
What does \mw{भेड़} mean? (\textquotedblleft{}A sheep\textquotedblright{}.) Say it once more.
\end{tcolorbox}
\section[\texorpdfstring{\mw{चिड़ी} (chiṛī)}{चिड़ी (chiṛī)}]{\mw{चिड़ी} (chiṛī) — a sparrow, a small bird}

\label{lesson:MW-C71-chiri}



\begin{quote}
Before the new one: say the Marwadi for a goat, then the Marwadi for a sheep.
\end{quote}

\subsection*{You'll want to know: \mw{चिड़ी}}
\textbf{\mw{चिड़ी}} — \emph{chiṛī} — \textquotedblleft{}a sparrow, a small bird\textquotedblright{}.

\begin{cousinweb}
\textbf{\mw{चिड़ी}} (\emph{chiṛī}) is a tadbhava, a word inherited from Sanskrit and reshaped by sound change. It comes from \emph{caṭaka}, \textquotedblleft{}sparrow\textquotedblright{}. CDIAL gives no Marwari form in this entry (CDIAL 4571); the match is through Hindi \emph{ciṛī}. CDIAL calls the word imitative and lists \emph{ciṛī} under a variant it reconstructs, \emph{*ciṭaka}.
\end{cousinweb}

\begin{grammarlens}[title={What you've built}]
One more word for this chapter.
\end{grammarlens}

\subsection*{Your turn}
\begin{itemize}
\raggedright
  \item \emph{Say it:} \emph{chiṛī}
  \item \emph{Say it:} \emph{chiṛī}, once more
  \item \emph{Say it:} say \emph{chiṛī}
  \item \emph{Recall:} say the Marwadi for a goat, then the Marwadi for a sheep, then say \emph{chiṛī} again
\end{itemize}

\begin{tcolorbox}[breakable,colback=teal!4,colframe=teal!35!black,title={Before you move on}]
What does \mw{चिड़ी} mean? (\textquotedblleft{}A sparrow, a small bird\textquotedblright{}.) Say it once more.
\end{tcolorbox}
\section[\texorpdfstring{\mw{नाहर} (nāhar)}{नाहर (nāhar)}]{\mw{नाहर} (nāhar) — a lion, a tiger}

\label{lesson:MW-C71-nahar}



\begin{quote}
Before the new one: say the Marwadi for a sheep, then the Marwadi for a sparrow.
\end{quote}

\subsection*{You'll want to know: \mw{नाहर}}
\textbf{\mw{नाहर}} — \emph{nāhar} — \textquotedblleft{}a lion, a tiger\textquotedblright{}.

\begin{cousinweb}
CDIAL lists Hindi \emph{nāhar}, \textquotedblleft{}tiger\textquotedblright{}, under two reconstructions, \emph{*nakhadara}, \textquotedblleft{}tearing with the claws\textquotedblright{} (CDIAL 6919), and \emph{*nakharin}, \textquotedblleft{}having claws\textquotedblright{} (CDIAL 6921). Both are built on Sanskrit \emph{nakha}, \textquotedblleft{}claw\textquotedblright{}; which one is the source is not settled, and CDIAL gives no Marwari form.
\end{cousinweb}

\begin{grammarlens}[title={What you've built}]
One more word for this chapter.
\end{grammarlens}

\subsection*{Your turn}
\begin{itemize}
\raggedright
  \item \emph{Say it:} \emph{nāhar}
  \item \emph{Say it:} \emph{nāhar}, once more
  \item \emph{Say it:} say \emph{nāhar}
  \item \emph{Recall:} say the Marwadi for a sheep, then the Marwadi for a sparrow, then say \emph{nāhar} again
\end{itemize}

\begin{tcolorbox}[breakable,colback=teal!4,colframe=teal!35!black,title={Before you move on}]
What does \mw{नाहर} mean? (\textquotedblleft{}A lion, a tiger\textquotedblright{}.) Say it once more.
\end{tcolorbox}
\section[\texorpdfstring{\mw{हाथी} (hāthī)}{हाथी (hāthī)}]{\mw{हाथी} (hāthī) — an elephant}

\label{lesson:MW-C71-hathi}



\begin{quote}
Before the new one: say the Marwadi for a sparrow, then the Marwadi for a lion.
\end{quote}

\subsection*{You'll want to know: \mw{हाथी}}
\textbf{\mw{हाथी}} — \emph{hāthī} — \textquotedblleft{}an elephant\textquotedblright{}.

\begin{cousinweb}
\textbf{\mw{हाथी}} (\emph{hāthī}) is a tadbhava, a word inherited from Sanskrit and reshaped by sound change. It comes from \emph{hastin}, \textquotedblleft{}having hands; elephant\textquotedblright{}. CDIAL lists Marwari \emph{hāthī} under it (CDIAL 14039). CDIAL asks whether the Sanskrit word reshaped an older \emph{karin}, perhaps of non-Aryan origin, tied to \emph{hasta}, \textquotedblleft{}hand\textquotedblright{}, by popular etymology.
\end{cousinweb}

\begin{grammarlens}[title={What you've built}]
One more word for this chapter.
\end{grammarlens}

\subsection*{Your turn}
\begin{itemize}
\raggedright
  \item \emph{Say it:} \emph{hāthī}
  \item \emph{Say it:} \emph{hāthī}, once more
  \item \emph{Say it:} say \emph{hāthī}
  \item \emph{Recall:} say the Marwadi for a sparrow, then the Marwadi for a lion, then say \emph{hāthī} again
\end{itemize}

\begin{tcolorbox}[breakable,colback=teal!4,colframe=teal!35!black,title={Before you move on}]
What does \mw{हाथी} mean? (\textquotedblleft{}An elephant\textquotedblright{}.) Say it once more.
\end{tcolorbox}
